This study began with a practical failure of listening. A feature film was playing through a small bass practice amplifier connected to the headphone socket of a computer, and for long stretches the only workable response to its soundtrack was to turn the volume all the way down and wait for the scene to end. Dialogue, when it returned, was too quiet to follow at the setting that had made the preceding gunfire tolerable. The volume went up for speech and down again for action, over and over, and some action sequences outlasted any patience for riding the control.

That experience is the only direct observation of listener behaviour in the monograph, and it is treated accordingly. It is reported once, in Chapter~\ref{ch:case}, as the report of a single listener on a single viewing through an unmeasured playback chain. Everything else in the book either measures the soundtrack itself, describes the instruments used to do so, or specifies what would have to be collected for the report to become evidence rather than motivation.

The book is organised as a sequence of instruments and their limits. Part~I states the problem and the conceptual framework. Part~II describes the analysis pipeline and, at some length, the ways its earlier versions produced confident nonsense: a swell detector that was an inverted transient detector, harshness scores that peaked in digital silence, and summary statistics whose values were fixed in advance by the thresholds that produced them. These failures are kept in the text because they are instructive and because a reader deserves to know how the current numbers were reached. Part~III applies the corrected instruments to the case. Part~IV specifies a response measure and an intervention, both of which were first sketched as speculative tools in July 2025 and are here developed into a research design. Part~V grades what has and has not been shown.

Numerical results that depend on the dialogue-referenced analysis of Chapter~\ref{ch:framework} are generated by \texttt{code/occupation.py} and \texttt{code/occupation-extras.py} and inserted automatically. \ifresults{In this build they have been generated from the analysis tables of the film.}{In this build they have not yet been generated, and each such value prints as \pending. Appendix~\ref{app:repro} gives the commands that fill them in.}

The code, the figures, and the manuscript source are released into the public domain.
